\documentclass[tikz,border=8pt]{standalone}
\usepackage{tikz}
\usepackage{amsmath}
\usepackage{pgfplots}
\pgfplotsset{compat=1.18}
\usepackage{ctex}
\usetikzlibrary{calc, positioning, arrows.meta, shapes, fit, backgrounds}

\begin{document}
\begin{tikzpicture}[
  every node/.style={font=\small},
  inode/.style={circle, draw=gray!60, fill=gray!15, minimum size=11mm, font=\small},
  lnode/.style={rectangle, draw=blue!70!black, fill=blue!10, minimum width=14mm,
                minimum height=12mm, rounded corners=2pt, font=\small\bfseries,
                align=center},
  >=Stealth
]

% =====================================================================
%  霍夫曼树结构
%  叶节点: A(0.40) B(0.20) C(0.15) D(0.15) E(0.10)
%  合并:
%    Step1: E(0.10) + D(0.15) = N1(0.25)
%    Step2: C(0.15) + B(0.20) = N2(0.35)
%    Step3: N1(0.25) + N2(0.35) = N3(0.60)
%    Step4: A(0.40) + N3(0.60) = root(1.00)
%  编码: A=0, E=100, D=101, C=110, B=111
% =====================================================================

% 根节点
\node[inode] (root) at (4.5, 8.0) {\small 1.00};

% 层1
\node[lnode] (A) at (1.5, 5.5) {A\\{\scriptsize 0.40}};
\node[inode] (N3) at (7.0, 5.5) {\small 0.60};

% 层2
\node[inode] (N1) at (5.5, 3.2) {\small 0.25};
\node[inode] (N2) at (8.5, 3.2) {\small 0.35};

% 层3（叶节点）
\node[lnode] (E) at (4.5, 1.0) {E\\{\scriptsize 0.10}};
\node[lnode] (D) at (6.5, 1.0) {D\\{\scriptsize 0.15}};
\node[lnode] (C) at (7.5, 1.0) {C\\{\scriptsize 0.15}};
\node[lnode] (B) at (9.5, 1.0) {B\\{\scriptsize 0.20}};

% 边 + 0/1 标注
\draw[->, thick] (root) -- node[midway, left, font=\scriptsize\bfseries]  {0} (A);
\draw[->, thick] (root) -- node[midway, right, font=\scriptsize\bfseries] {1} (N3);
\draw[->, thick] (N3)   -- node[midway, left, font=\scriptsize\bfseries]  {0} (N1);
\draw[->, thick] (N3)   -- node[midway, right, font=\scriptsize\bfseries] {1} (N2);
\draw[->, thick] (N1)   -- node[midway, left, font=\scriptsize\bfseries]  {0} (E);
\draw[->, thick] (N1)   -- node[midway, right, font=\scriptsize\bfseries] {1} (D);
\draw[->, thick] (N2)   -- node[midway, left, font=\scriptsize\bfseries]  {0} (C);
\draw[->, thick] (N2)   -- node[midway, right, font=\scriptsize\bfseries] {1} (B);

% 整体标题
\node[font=\large\bfseries] at (5.5, 9.5) {霍夫曼编码树};

% =====================================================================
% 右侧编码说明（纯节点，不用 tabular）
% =====================================================================
\node[font=\normalsize\bfseries] at (14.5, 8.5) {编码表};

\node[font=\small, align=left,
      fill=white, draw=gray!50, thin, rounded corners=2pt, inner sep=8pt]
  at (14.5, 5.5) {
    符号 A：编码 \texttt{0}\quad（1 bit）\\[3pt]
    符号 E：编码 \texttt{100}\quad（3 bit）\\[3pt]
    符号 D：编码 \texttt{101}\quad（3 bit）\\[3pt]
    符号 C：编码 \texttt{110}\quad（3 bit）\\[3pt]
    符号 B：编码 \texttt{111}\quad（3 bit）
  };

\node[font=\small,
      fill=white, draw=gray!50, thin, rounded corners=2pt, inner sep=6pt]
  at (14.5, 2.0) {
    平均码长 $\bar{L}=0.4\times1+0.6\times3=2.2$ bit
  };

% 底部公式说明
\node[font=\small,
      fill=white, draw=gray!50, thin, rounded corners=2pt, inner sep=5pt]
  at (7.0, -0.5) {
    贪心合并：每次选取概率最小的两个节点合并，保证前缀码最优
  };

\end{tikzpicture}
\end{document}
